% !TeX program = xelatex
\documentclass[11pt]{article}
\usepackage{geometry}
\geometry{a4paper, margin=2.6cm}

\usepackage{fontspec}
\usepackage{unicode-math}
\setmainfont{TeX Gyre Pagella}
\setmathfont{TeX Gyre Pagella Math}
\newfontfamily\isafont{DejaVu Sans Mono}[Scale=MatchLowercase]
\newfontfamily\isafallbackA{DejaVu Sans}[Scale=MatchLowercase]
\newfontfamily\isafallbackB{TeX Gyre Pagella Math}[Scale=MatchLowercase]
\usepackage{newunicodechar}
\newcommand{\gl}[1]{\iffontchar\font#1\relax\char#1\relax\else
  {\isafallbackA\iffontchar\font#1\relax\char#1\relax\else
   {\isafallbackB\char#1\relax}\fi}\fi}
\newunicodechar{⟹}{\gl{"27F9}}
\newunicodechar{⟸}{\gl{"27F8}}
\newunicodechar{⟶}{\gl{"27F6}}
\newunicodechar{⟷}{\gl{"27F7}}
\newunicodechar{⇒}{\gl{"21D2}}
\newunicodechar{⟦}{\gl{"27E6}}
\newunicodechar{⟧}{\gl{"27E7}}
\newunicodechar{⟨}{\gl{"27E8}}
\newunicodechar{⟩}{\gl{"27E9}}
\newunicodechar{‹}{\gl{"2039}}
\newunicodechar{›}{\gl{"203A}}
\newunicodechar{⋀}{\gl{"22C0}}
\newunicodechar{⊩}{\gl{"22A9}}
\newunicodechar{⊢}{\gl{"22A2}}
\newunicodechar{▹}{\gl{"25B9}}
\newunicodechar{↦}{\gl{"21A6}}
\newunicodechar{⇓}{\gl{"21D3}}
\newunicodechar{∅}{\gl{"2205}}
\newunicodechar{∈}{\gl{"2208}}
\newunicodechar{≡}{\gl{"2261}}
\newunicodechar{≠}{\gl{"2260}}
\newunicodechar{≤}{\gl{"2264}}
\newunicodechar{≥}{\gl{"2265}}
\newunicodechar{∧}{\gl{"2227}}
\newunicodechar{∨}{\gl{"2228}}
\newunicodechar{¬}{\gl{"00AC}}
\newunicodechar{∀}{\gl{"2200}}
\newunicodechar{∃}{\gl{"2203}}
\newunicodechar{λ}{\gl{"03BB}}

\usepackage{amsmath, amsthm}
\usepackage{xcolor}
\usepackage{fvextra}
\usepackage[breakable,skins]{tcolorbox}
\usepackage{booktabs}
\usepackage{array}
\usepackage{microtype}
\usepackage[hidelinks]{hyperref}

\definecolor{isabg}{gray}{0.955}
\definecolor{isarule}{gray}{0.80}
\definecolor{thmrule}{RGB}{20,50,110}

\fvset{formatcom=\isafont, fontsize=\footnotesize,
       breaklines=true, breakanywhere=true,
       breaksymbolleft=\textcolor{isarule}{\tiny\ensuremath{\hookrightarrow}}}

\newtcolorbox{isabox}{
  breakable, enhanced, sharp corners,
  colback=isabg, colframe=isarule,
  boxrule=0pt, leftrule=1.2pt,
  left=6pt, right=2pt, top=3pt, bottom=3pt,
  before skip=1.0em, after skip=1.0em,
}

\tolerance=1500
\emergencystretch=2.5em

\DeclareRobustCommand{\isa}[1]{\begingroup\isafont\small
  \renewcommand{\_}{\textunderscore\allowbreak}#1\endgroup}
\newcommand{\thy}[1]{%
  \begin{isabox}\VerbatimInput{snippets/#1.thy}\end{isabox}}

\theoremstyle{definition}
\newtheorem{definitionx}{Definition}[section]
\newtheorem{remarkx}[definitionx]{Remark}
\theoremstyle{plain}
\newtheorem{theoremx}[definitionx]{Theorem}
\newtheorem{lemmax}[definitionx]{Lemma}

\newcommand{\irule}[3]{%
  \ensuremath{\dfrac{\rule{0pt}{2.1ex}#1}{\rule{0pt}{2.4ex}#2}}%
  \;\text{\scriptsize(\textsc{#3})}}
\newcommand{\iaxiom}[2]{%
  \ensuremath{\dfrac{\rule{0pt}{2.1ex}\strut}{\rule{0pt}{2.4ex}#1}}%
  \;\text{\scriptsize(\textsc{#2})}}
\newcommand{\dis}[2]{%
  \ensuremath{\begin{array}{c}[#1]\\[-1pt]\vdots\\[-1pt]#2\end{array}}}
\newcommand{\sep}{\quad}
\newenvironment{rules}%
  {\begin{center}\begin{tabular}{@{}c@{\hspace{2.2em}}c@{}}}%
  {\end{tabular}\end{center}}
\newenvironment{rulesiii}%
  {\begin{center}\begin{tabular}{@{}c@{\hspace{1.6em}}c@{\hspace{1.6em}}c@{}}}%
  {\end{tabular}\end{center}}

\newcommand{\N}{\ensuremath{\,\mathsf{N}}}
\newcommand{\B}{\ensuremath{\,\mathsf{B}}}
\newcommand{\Suc}[1]{\ensuremath{\mathsf{S}(#1)}}
\newcommand{\Pre}[1]{\ensuremath{\mathsf{P}(#1)}}
\newcommand{\seq}{\ensuremath{\vdash}}
\newcommand{\ent}{\ensuremath{\Rightarrow}}
\newcommand{\enc}[1]{\ensuremath{\ulcorner #1 \urcorner}}
\newcommand{\ite}[3]{\ensuremath{\mathsf{if}\;#1\;\mathsf{then}\;#2\;\mathsf{else}\;#3}}
\newcommand{\fn}[1]{\ensuremath{\mathsf{#1}}}
\newcommand{\dfeq}{\ensuremath{\;\triangleq\;}}
\newcommand{\rec}{\ensuremath{\;:=\;}}
\newcommand{\GA}{\textsf{GA}}
\newcommand{\QGA}{\textsf{QGA}}
\newcommand{\BGA}{\textsf{BGA}}
\newcommand{\PGA}{\textsf{PGA}}
\newcommand{\GD}{\textsf{GD}}


\title{\bfseries Self-Consistency of Basic Grounded Arithmetic\\[2pt]
       \large A consistency proof for \BGA{} carried out inside grounded
       arithmetic in Isabelle/Pure}
\author{Omar Muhammad}
\date{\today}

\begin{document}
\maketitle

\begin{abstract}
\noindent
Grounded arithmetic admits arbitrary recursive definitions by weakening
inference rules with \emph{habeas quid} termination premises.  The existing
mechanized consistency results for grounded arithmetic are carried out in
Isabelle/HOL, that is, in a classical metatheory far stronger than the object
logic.  This report documents a development that carries the argument out
inside grounded arithmetic itself.  \QGA, an axiomatization of grounded
arithmetic with object-level entailment and quantifiers, is embedded into
Isabelle/Pure in \texttt{GD.thy}.  Basic grounded arithmetic \BGA{} is
arithmetized on top of it in \texttt{BGA\_on\_GA.thy}: terms, formulas,
judgments, proofs, the operational semantics and the proof checker are encoded
as \QGA{} natural numbers, and the checker is proved sound with respect to the
encoded semantics under a small set of abstract assumptions about the encoding.
A concrete Cantor-pairing encoding in \texttt{encode.thy} discharges those
assumptions and yields
\[
  \neg\bigl(\textsf{is\_valid\_proof}\;p_1\;\langle \textsf{Nil},\,
  \enc{a = b}\rangle \;\wedge\;
  \textsf{is\_valid\_proof}\;p_2\;\langle \textsf{Nil},\, \enc{a \neq b}\rangle\bigr)
\]
as a theorem of \QGA.  Sections \ref{sec:qga} to \ref{sec:encode} state the
rules and definitions of each layer and outline the consistency proof.  This is
a first instalment: the detailed proofs, and the reduction of \QGA{} to \BGA{}
that closes the self-consistency argument, are still to be written up.
\end{abstract}

\tableofcontents

\section{Overview}
\label{sec:overview}

\subsection{The three layers}

The development consists of three theories, each importing the previous one.

\begin{center}
\begin{tabular}{@{}lll@{}}
\toprule
\textbf{Theory} & \textbf{Imports} & \textbf{Contributes}\\
\midrule
\texttt{GD.thy}          & \texttt{Pure}        & \QGA, an embedding of grounded arithmetic\\
                         &                      & in Isabelle/Pure\\[2pt]
\texttt{BGA\_on\_GA.thy} & \texttt{GD}          & \BGA{} arithmetized in \QGA, and the\\
                         &                      & soundness proof for its proof checker\\[2pt]
\texttt{encode.thy}      & \texttt{BGA\_on\_GA} & a concrete encoding discharging every\\
                         &                      & assumption, and the theorem\\
\bottomrule
\end{tabular}
\end{center}

\noindent
\texttt{BGA\_on\_GA.thy} does the heavy lifting.  It never commits to a
particular G\"odel numbering.  It fixes an interface (\isa{tag\_T},
\isa{load\_T}, \isa{pack\_T} and their formula counterparts) together with the
smallest set of assumptions under which the soundness argument goes through,
and proves everything relative to that interface.  \texttt{encode.thy} then
pays the arithmetic cost once, by exhibiting one encoding that satisfies the
assumptions, and instantiates the result.

\subsection{What is proved}

The final theorem, \isa{BGA\_syntactically\_consistent}, says that there is no
pair of numbers $p_1$, $p_2$ such that $p_1$ codes a \BGA{} proof of $a = b$
from no hypotheses and $p_2$ codes a \BGA{} proof of $a \neq b$ from no
hypotheses.  It is a statement about the encoded proof checker, and it is
proved in \QGA: every quantifier in it is a \QGA{} quantifier, every function
used to state it is a \QGA{} recursive definition, and every step of the proof
is an application of a \QGA{} rule as axiomatized in \texttt{GD.thy}.
Isabelle/Pure supplies the framework in which those rules are stated and
applied, and contributes no logical strength of its own.

\QGA{} is stronger than \BGA, so this is a consistency proof of a weaker system
inside a stronger one, in the ordinary G\"odelian arrangement.  Its interest
lies elsewhere: it shows that grounded arithmetic is expressive enough to carry
a full metamathematical argument about a formal system, with no classical
ambient logic available at any point.  The step that would turn it into a
self-consistency result, a reduction of \QGA{} back to \BGA, is
\S\ref{sec:reduction} and is left for later work.

\subsection{Reading the rules}
\label{sec:levels}

Isabelle/Pure is the framework in which \QGA{} is stated, and three of its
symbols appear throughout.  \isa{A ⟹ B} is framework implication, and it is how
a \QGA{} rule ``from $A$ infer $B$'' is written, with hypothesis contexts left
implicit because Pure manages them.  \isa{⋀x. A} binds a variable in a rule
schema, and is how an eigenvariable side condition is expressed.  \isa{a ≡ b}
is the definitional equality of the framework, used when a \QGA{} notion is
abbreviated rather than axiomatized.  None of the three is a connective of
\QGA.

\QGA{} has its own logical vocabulary, introduced in \S\ref{sec:qga}: the
primitives \isa{∨} and \isa{¬}, the connectives \isa{∧}, \isa{⟶} and \isa{⟷}
defined from them, the quantifiers \isa{∀} and \isa{∃}, and an entailment
constant \isa{⊢} which is a shorthand for framework implication written as a
formula.

In the natural-deduction presentations below, \QGA{} rules are written with
premises above the line and the conclusion below it and with no explicit
context.  A premise discharged by a rule is written in brackets above the
derivation that uses it.  For \BGA, whose judgments carry an explicit list of
hypotheses, rules are written as sequents $\Gamma \seq \varphi$.

All Isabelle snippets are extracted verbatim from the theories by
\texttt{extract.sh}, which replaces Isabelle's ASCII escapes by the
corresponding Unicode glyphs and changes nothing else.

\section{\QGA}
\label{sec:qga}

\texttt{GD.thy} axiomatizes grounded arithmetic as an object logic over
Isabelle/Pure.  Beyond the propositional and arithmetic core it adds
first-order quantifiers and a shorthand \isa{⊢} for framework implication, and
we call the resulting system \QGA.

\subsection{Formulas and terms}

Two Pure types carry the embedding.  \isa{o} is the type of \QGA{} formulas and
\isa{num} is the type of \QGA{} terms.  A \isa{judgment} declaration supplies
the coercion \isa{Trueprop :: o ⇒ prop} that lets a formula be asserted at the
framework level, so that a \QGA{} formula can appear as the conclusion of a
Pure rule.

\subsection{Propositional rules}

Disjunction and negation are primitive.

\begin{rulesiii}
\irule{P}{P \vee Q}{$\vee$I1} &
\irule{Q}{P \vee Q}{$\vee$I2} &
\irule{\neg P \sep \neg Q}{\neg (P \vee Q)}{$\vee$I3}\\[3.2ex]
\irule{\neg(P \vee Q)}{\neg P}{$\vee$E2} &
\irule{\neg(P \vee Q)}{\neg Q}{$\vee$E3} &
\irule{P \vee Q \sep \dis{P}{R} \sep \dis{Q}{R}}{R}{$\vee$E1}\\[4.6ex]
\irule{P}{\neg\neg P}{$\neg\neg$I} &
\irule{\neg\neg P}{P}{$\neg\neg$E} &
\irule{P \sep \neg P}{Q}{exF}
\end{rulesiii}

\thy{gd-prop}

\noindent
The introduction rules are the classical ones and the De Morgan directions
$\vee$E2 and $\vee$E3 are available unconditionally.  What is weakened is case
analysis.  $\vee$E1 needs an actual proof of $P \vee Q$, and $P \vee Q$ is not
derivable for arbitrary $P$: the law of excluded middle is a judgment to be
established rather than an axiom, and it is written \isa{P B}.

\subsection{Equality and the habeas quid judgments}

Equality is a binary predicate on terms, \isa{num ⇒ num ⇒ o}, axiomatized by
substitution and symmetry.

\begin{rules}
\irule{a = b \sep Q\,a}{Q\,b}{=E} &
\irule{a = b}{b = a}{=S}
\end{rules}

\thy{gd-eq}

\noindent
Reflexivity is absent, and this is the pivot of the whole system.
Transitivity follows from one substitution, \isa{eq\_trans: a = b ⟹ b = c ⟹ a =
c}.  The third axiom, \isa{eq\_reflection}, lifts a derivable equation to the
framework equality \isa{≡}, which is what lets Isabelle's simplifier operate on
\QGA{} goals; it is a statement about the framework and adds nothing to the
object logic.

The two habeas quid judgments and the remaining connectives are definitions.

\thy{gd-defs}

\begin{definitionx}[Habeas quid]
\isa{a N}, defined as \isa{a = a}, asserts that the term \isa{a} denotes a
natural number, equivalently that the computation it names terminates.
\isa{p B}, defined as \isa{p ∨ ¬p}, asserts that the formula \isa{p} is
grounded.
\end{definitionx}

\noindent
Since \isa{a N} unfolds to \isa{a = a}, reflexivity of equality says exactly
that every term terminates, which is why it cannot be an axiom of a logic with
unrestricted recursion.  Conjunction, implication and the biconditional are
defined from $\vee$ and $\neg$ in the classical way, and their rules are then
derived.  Conjunction comes out unconditional, while implication and the
biconditional acquire habeas quid premises, because their derivations pass
through $\vee$E1.

\begin{rulesiii}
\irule{p \sep q}{p \wedge q}{$\wedge$I} &
\irule{p \wedge q}{p}{$\wedge$E1} &
\irule{p \wedge q}{q}{$\wedge$E2}\\[3.2ex]
\irule{a \B \sep \dis{a}{b}}{a \longrightarrow b}{$\longrightarrow$I} &
\irule{a \longrightarrow b \sep a}{b}{$\longrightarrow$E} &
\irule{a \B \sep b \B \sep \dis{a}{b} \sep \dis{b}{a}}
      {a \longleftrightarrow b}{$\longleftrightarrow$I}
\end{rulesiii}

\noindent
Two further derived rules are used constantly.  Case analysis on a grounded
formula,
\begin{center}
\isa{cases\_bool: ⟦q B; q ⟹ p; ¬q ⟹ p⟧ ⟹ p},
\end{center}
and grounded proof by contradiction,
\begin{center}
\isa{grounded\_contradiction: ⟦p B; ¬p ⟹ q; ¬p ⟹ ¬q⟧ ⟹ p},\qquad
\isa{contradiction: ⟦p B; ¬p ⟹ False⟧ ⟹ p}.
\end{center}
Both require the formula in question to be grounded.  The final consistency
argument of \S\ref{sec:consistency} is an application of \isa{contradiction},
and the groundedness premise it needs is the reason the abstract development
assumes that proof checking is boolean.

\subsection{Natural numbers}

\begin{rulesiii}
\iaxiom{0 \N}{0I} &
\irule{\Suc{a} = \Suc{b}}{a = b}{S-inj} &
\irule{a = b}{\Suc{a} = \Suc{b}}{S-cong}\\[3.2ex]
\irule{a = b}{\Pre{a} = \Pre{b}}{P-cong} &
\irule{a \N \sep b \N}{(a = b) \B}{eqBool} &
\irule{a \N}{\Suc{a} \neq 0}{S$\neq$0}\\[3.2ex]
\irule{a \N}{\Pre{\Suc{a}} = a}{PS} &
\iaxiom{\Pre{0} = 0}{P0} &
\irule{(a = b) \B}{a \N \wedge b \N}{eqE}\\[3.2ex]
\multicolumn{1}{c}{\irule{\Pre{a} \N}{a \N}{predTI}} &
\multicolumn{2}{c}{\irule{a \N \sep Q\,0 \sep \dis{x \N,\; Q\,x}{Q\,\Suc{x}}}{Q\,a}{Ind}}
\end{rulesiii}

\thy{gd-nat}

\noindent
\isa{eqBool} is the rule that converts termination into decidability: two
terminating terms are either equal or not.  \isa{eqE} and \isa{predTIE} run in
the opposite direction, recovering termination from groundedness of an equation
and from termination of a predecessor.  Induction carries a habeas quid premise
on the subject of the induction, and its step case may assume both \isa{x N}
and the induction hypothesis.

Truth and falsity are definitions.

\thy{gd-truefalse}

\noindent
Derived facts of constant use are \isa{zeroRefl: zero = zero},
\isa{natS: a N ⟹ S a N}, \isa{natSI: S a N ⟹ a N}, \isa{natP: a N ⟹ P a N},
and the pair \isa{eq\_impl\_term: a = b ⟹ a N} and
\isa{eq\_impl\_term2: a = b ⟹ b N}, by which a derivable equation witnesses the
termination of both of its sides.

\subsection{Entailment}
\label{sec:entails}

Entailment is a shorthand: its two rules make \isa{a ⊢ b} and the framework
implication \isa{a ⟹ b} interderivable, so it is framework implication written
as a formula, available wherever a formula is required.

\begin{rules}
\irule{\dis{a}{b}}{a \seq b}{$\seq$I} &
\irule{a \seq b \sep a}{b}{$\seq$E}
\end{rules}

\thy{gd-entails}

\noindent
Unlike \isa{⟶} it has no groundedness premise on its introduction rule, which
is what makes it usable in \S\ref{sec:list}, where the antecedent of the
implication being inducted on cannot be shown grounded.

\subsection{Quantifiers}

\begin{rules}
\irule{\dis{x \N}{Q\,x}}{\forall x.\,Q\,x}{$\forall$I} &
\irule{\forall x.\,Q\,x \sep a \N}{Q\,a}{$\forall$E}\\[4.6ex]
\irule{a \N \sep Q\,a}{\exists x.\,Q\,x}{$\exists$I} &
\irule{\exists x.\,Q\,x \sep \dis{a \N,\; Q\,a}{R}}{R}{$\exists$E}
\end{rules}

\thy{gd-quant}

\noindent
Both quantifiers range over \isa{num} and are relativized to terminating
witnesses.  $\forall$I proves the body for an arbitrary terminating
eigenvariable, $\forall$E instantiates only at terminating terms, and
$\exists$I must exhibit a terminating witness.  The negation-exchange rules for
the quantifiers appear in the source only as commented-out axioms.  They are
not part of the axiomatization and nothing downstream uses them.

\subsection{Conditionals}

The if-zero-else construct is written \isa{if c then a else b} and is
polymorphic in the branches, so it serves both for term-valued and for
formula-valued conditionals.

\begin{rules}
\irule{c \sep a \N}{(\text{if } c \text{ then } a \text{ else } b) = a}{condI1} &
\irule{\neg c \sep b \N}{(\text{if } c \text{ then } a \text{ else } b) = b}{condI2}\\[3.2ex]
\irule{c \sep (\text{if } c \text{ then } a \text{ else } b) \N}{a \N}{condE1} &
\irule{(\text{if } c \text{ then } a \text{ else } b) \N}{c \B}{condE3}\\[3.2ex]
\irule{c \sep Q\,(\text{if } c \text{ then } a \text{ else } b)}{Q\,a}{cond-then-E} &
\irule{c \sep Q\,a}{Q\,(\text{if } c \text{ then } a \text{ else } b)}{cond-then-I}
\end{rules}

\thy{gd-cond}

\noindent
Each rule has a formula-level counterpart in which $=$ is replaced by
$\longleftrightarrow$ and $\mathsf{N}$ by $\mathsf{B}$, and each ``then'' rule
has an ``else'' twin.
The last pair is the lazy form: once the guard is decided, an arbitrary context
\isa{Q} transports across the conditional.  These are the rules that let a
recursive definition be unfolded one branch at a time without evaluating the
other branch.  The two totality lemmas
\begin{center}
\isa{condT: ⟦c B; a N; b N⟧ ⟹ (if c then a else b) N},\qquad
\isa{condTB: ⟦c B; a B; b B⟧ ⟹ (if c then a else b) B}
\end{center}
are derived rather than assumed, and they carry most of the groundedness proofs
about the encoded checkers of \S\ref{sec:bga}.

\subsection{The definitional mechanism}
\label{sec:defmech}

A recursive definition is admitted with no termination side condition.  The
constant is introduced by \isa{axiomatization} and its defining equation is an
axiom of the form \isa{f x := body}.

\begin{rules}
\irule{a := b \sep Q\,b}{Q\,a}{defE} &
\irule{a := b \sep Q\,a}{Q\,b}{defI}
\end{rules}

\thy{gd-def}

\noindent
\isa{a := b} does not assert \isa{a = b}.  It asserts that \isa{a} and \isa{b}
are interchangeable in every context \isa{Q}, whether or not either denotes a
value.  That is what makes \isa{omega := omega} harmless: \isa{omega} is
interchangeable with itself and \isa{omega N} is underivable.  The two rules
are the unfolding and folding directions, and the \isa{unfold\_def} method of
\texttt{unfold\_def.ML} applies them.

\subsection{Arithmetic}

The standard functions are introduced by exactly this mechanism.

\thy{gd-arith}

\noindent
Comparison returns a term rather than a formula, so \isa{x ≤ y} evaluates to
\isa{1} or \isa{0} and the associated formula is \isa{x ≤ y = 1}.  The constant
\isa{omega} is included to exercise the claim that a divergent definition is
admissible.  Termination lemmas such as \isa{add\_terminates},
\isa{leq\_terminates} and \isa{less\_terminates} are proved by induction, as is
termination of Ackermann's function, which is defined in the same file and is
not primitive recursive.

\subsection{Induction principles}

Two derived induction principles carry most of the later work.

\thy{gd-strong-induction}
\thy{gd-list-induct}

\noindent
Strong induction is stated with the bounded hypothesis \isa{y ≤ x = 1}.  It is
what makes recursion over the structure of an encoded term legitimate: a
component of a code is numerically smaller than the code, so structural
recursion becomes bounded numerical recursion.  \isa{list\_induct} is derived
from it through the encoding of lists below.

\subsection{Pairing}
\label{sec:cpair}

Pairing is the bridge from numbers to syntax.  It is defined recursively rather
than by a closed formula.

\thy{gd-cpair}

\noindent
The syntax \isa{⟨x, y⟩} abbreviates \isa{cpair x y} and nests to the right for
tuples.  The inverses are recursive as well.

\thy{gd-cpxcpy}

\noindent
\isa{cpi n x} projects the $n$th component of a right-nested tuple by iterating
\isa{cpy} and applying \isa{cpx}.

\thy{gd-cpi}

\noindent
The properties established for the pairing are exactly the ones the encoding of
\S\ref{sec:encode} needs.

\begin{center}
\begin{tabular}{@{}ll@{}}
\toprule
\isa{cpair\_terminates} & \isa{x N ⟹ y N ⟹ ⟨x, y⟩ N}\\
\isa{cpair\_suc}        & \isa{x N ⟹ y N ⟹ ⟨x, S(y)⟩ = ⟨x, y⟩ + x + S(y) + 1}\\
\isa{cpx\_proj}, \isa{cpy\_proj} & the projection equations\\
\isa{cpair\_inj}        & injectivity\\
\isa{cpair\_surjective} & every number is a pair\\
\isa{cpair\_reconstr}   & \isa{a N ⟹ ⟨cpx a, cpy a⟩ = a}\\
\isa{cpx\_mono}, \isa{cpy\_mono} & each projection is bounded by its argument\\
\bottomrule
\end{tabular}
\end{center}

\thy{gd-cpair-proj}
\thy{gd-cpair-inj}
\thy{gd-cpair-surj}
\thy{gd-cp-mono}

\noindent
The Cantor pairing is a bijection between $\mathbb{N}^2$ and $\mathbb{N}$, and
both directions of that fact are used later:  injectivity gives the
decomposition properties of the encoding, and surjectivity gives
\isa{cpair\_reconstr}, which says that a number can always be taken apart and
put back together.  Monotonicity feeds strong induction: since \isa{cpx x ≤ x}
and \isa{cpy x ≤ x}, a recursion that descends into components terminates.

\subsection{Lists}

Hypothesis sets, assignments, definition lists and proofs are all lists, so
lists are the last piece of \QGA{} infrastructure.  \isa{Nil} is $0$ and
\isa{Cons} is a shifted pair, which makes the list encoding a bijection with
$\mathbb{N}$ and removes the need for any well-formedness predicate on list
codes.

\thy{gd-list}

\noindent
The derived facts are the expected ones: \isa{nil\_nat},
\isa{cons\_nat: ⟦n N; xs N⟧ ⟹ Cons n xs N}, disjointness of the two
constructors, the projection equations \isa{list\_hd (Cons n xs) = n} and
\isa{list\_tl (Cons n xs) = xs}, reconstruction
\isa{cons\_reconstr: ⟦x N; ¬(x = Nil)⟧ ⟹ Cons (list\_hd x) (list\_tl x) = x},
injectivity \isa{cons\_inj}, exhaustiveness
\isa{list\_cases: x N ⟹ (x = Nil) ∨ (∃n xs. x = Cons n xs)}, and
\isa{list\_induct}.  Together these are the characteristic properties of an
inductive datatype, obtained with no datatype package.

Four list functions are defined recursively and used throughout
\S\ref{sec:bga}.

\thy{gd-mem}
\thy{gd-nth}
\thy{gd-len}
\thy{gd-subset}

\noindent
The two habeas quid results about them that the proof checker relies on are
\isa{mem\_bool: ⟦x N; G N⟧ ⟹ mem x G B} and
\isa{nth\_in\_range\_N: ⟦xs N; i < len xs = 1⟧ ⟹ nth i xs N}.

\subsection{Summary}
\label{sec:qga-summary}

Table \ref{tab:qga-defs} collects the definitions of \texttt{GD.thy} and
Table \ref{tab:qga-rules} its inference rules, both in the notation of \QGA{}
rather than in Isabelle syntax.  A definition written with $\triangleq$ is an
abbreviation, unfolded by the framework.  A definition written with $:=$ is a
\QGA{} recursive definition in the sense of \S\ref{sec:defmech}, admitted with
no termination obligation and used only through \textsc{defE} and
\textsc{defI}.

\begin{table}[p]
\centering
\begin{tabular}{@{}ll@{}}
\toprule
\multicolumn{2}{@{}l}{\textbf{Abbreviations}}\\
\midrule
$a \neq b$          & $\triangleq \neg (a = b)$\\
$p \B$              & $\triangleq p \vee \neg p$\\
$a \N$              & $\triangleq a = a$\\
$p \wedge q$        & $\triangleq \neg(\neg p \vee \neg q)$\\
$p \longrightarrow q$   & $\triangleq \neg p \vee q$\\
$p \longleftrightarrow q$ & $\triangleq (p \longrightarrow q) \wedge (q \longrightarrow p)$\\
$\fn{True}$         & $\triangleq 0 = 0$\\
$\fn{False}$        & $\triangleq \Suc{0} = 0$\\
$x > y$             & $\triangleq 1 - (x \leq y)$\\
$x \geq y$          & $\triangleq 1 - (x < y)$\\
$\fn{Nil}$          & $\triangleq 0$\\
$\fn{Cons}\;n\;xs$  & $\triangleq \langle n, xs\rangle + 1$\\
$\fn{hd}\;x$        & $\triangleq \fn{cpx}\,\Pre{x}$\\
$\fn{tl}\;x$        & $\triangleq \fn{cpy}\,\Pre{x}$\\
\midrule
\multicolumn{2}{@{}l}{\textbf{Recursive definitions}}\\
\midrule
$x + y$   & $:= \ite{y = 0}{x}{\Suc{x + \Pre{y}}}$\\
$x - y$   & $:= \ite{y = 0}{x}{\Pre{x - \Pre{y}}}$\\
$x * y$   & $:= \ite{y = 0}{0}{x + (x * \Pre{y})}$\\
$x \leq y$ & $:= \ite{x = 0}{1}{\ite{y = 0}{0}{\Pre{x} \leq \Pre{y}}}$\\
$x < y$   & $:= \ite{y = 0}{0}{\ite{x = 0}{1}{\Pre{x} < \Pre{y}}}$\\
$\fn{div}\;x\;y$ & $:= \ite{(x < y) = 1}{0}{\Suc{\fn{div}\;(x - y)\;y}}$\\
$\omega$  & $:= \omega$\\
$\langle x, y\rangle$ & $:= \ite{y = 0}{\fn{div}\;(x * \Suc{x})\;2}
                            {\langle x, \Pre{y}\rangle + x + y + 1}$\\
$\fn{cpx}\;x$ & $:= \ite{x = 0}{0}{\ite{\fn{cpx}\,\Pre{x} = 0}
                    {\Suc{\fn{cpy}\,\Pre{x}}}{\Pre{\fn{cpx}\,\Pre{x}}}}$\\
$\fn{cpy}\;x$ & $:= \ite{\fn{cpx}\,\Pre{x} = 0}{0}{\Suc{\fn{cpy}\,\Pre{x}}}$\\
$\fn{cpi}'\;n\;x$ & $:= \ite{n = 0}{0}{\ite{n = 1}{x}{\fn{cpy}\,(\fn{cpi}'\;(n-1)\;x)}}$\\
$x \in G$ & $:= \ite{G = \fn{Nil}}{\fn{False}}
                {\ite{\fn{hd}\,G = x}{\fn{True}}{x \in \fn{tl}\,G}}$\\
$\fn{nth}\;i\;xs$ & $:= \ite{xs = \fn{Nil}}{0}
                {\ite{i = 0}{\fn{hd}\,xs}{\fn{nth}\;(i-1)\;(\fn{tl}\,xs)}}$\\
$\fn{len}\;xs$ & $:= \ite{xs = \fn{Nil}}{0}{\Suc{\fn{len}\,(\fn{tl}\,xs)}}$\\
$\fn{subset}\;G'\;G$ & $:= \ite{G' = \fn{Nil}}{\fn{True}}
                {\fn{hd}\,G' \in G \wedge \fn{subset}\;(\fn{tl}\,G')\;G}$\\
$\fn{ack}\;x\;y$ & $:= \ite{x = 0}{y + 1}{\ite{y = 0}{\fn{ack}\,\Pre{x}\,1}
                {\fn{ack}\,\Pre{x}\,(\fn{ack}\;x\;\Pre{y})}}$\\
\bottomrule
\end{tabular}
\caption{Definitions of \QGA.  $\langle x, y\rangle$ is the Cantor pairing and
$\fn{cpx}$, $\fn{cpy}$ its projections.}
\label{tab:qga-defs}
\end{table}

\begin{table}[p]
\centering\small
\begin{tabular}{@{}c@{\hspace{1.4em}}c@{\hspace{1.4em}}c@{}}
\multicolumn{3}{@{}l}{\textbf{Propositional}}\\[0.9ex]
\irule{P}{P \vee Q}{$\vee$I1} &
\irule{Q}{P \vee Q}{$\vee$I2} &
\irule{\neg P \sep \neg Q}{\neg (P \vee Q)}{$\vee$I3}\\[2.4ex]
\irule{\neg(P \vee Q)}{\neg P}{$\vee$E2} &
\irule{\neg(P \vee Q)}{\neg Q}{$\vee$E3} &
\irule{P \vee Q \sep \dis{P}{R} \sep \dis{Q}{R}}{R}{$\vee$E1}\\[3.6ex]
\irule{P}{\neg\neg P}{$\neg\neg$I} &
\irule{\neg\neg P}{P}{$\neg\neg$E} &
\irule{P \sep \neg P}{Q}{exF}\\[2.8ex]
\multicolumn{3}{@{}l}{\textbf{Equality}}\\[0.9ex]
\irule{a = b \sep Q\,a}{Q\,b}{=E} &
\irule{a = b}{b = a}{=S} &
\\[2.8ex]
\multicolumn{3}{@{}l}{\textbf{Natural numbers}}\\[0.9ex]
\iaxiom{0 \N}{0I} &
\irule{\Suc{a} = \Suc{b}}{a = b}{S-inj} &
\irule{a = b}{\Suc{a} = \Suc{b}}{S-cong}\\[2.4ex]
\irule{a = b}{\Pre{a} = \Pre{b}}{P-cong} &
\irule{a \N \sep b \N}{(a = b) \B}{eqBool} &
\irule{a \N}{\Suc{a} \neq 0}{S$\neq$0}\\[2.4ex]
\irule{a \N}{\Pre{\Suc{a}} = a}{PS} &
\iaxiom{\Pre{0} = 0}{P0} &
\irule{(a = b) \B}{a \N \wedge b \N}{eqE}\\[2.4ex]
\irule{\Pre{a} \N}{a \N}{predTI} &
\multicolumn{2}{c}{\irule{a \N \sep Q\,0 \sep \dis{x \N,\; Q\,x}{Q\,\Suc{x}}}{Q\,a}{Ind}}\\[2.8ex]
\multicolumn{3}{@{}l}{\textbf{Entailment and quantifiers}}\\[0.9ex]
\irule{\dis{a}{b}}{a \seq b}{$\seq$I} &
\irule{a \seq b \sep a}{b}{$\seq$E} &
\irule{\dis{x \N}{Q\,x}}{\forall x.\,Q\,x}{$\forall$I}\\[3.6ex]
\irule{\forall x.\,Q\,x \sep a \N}{Q\,a}{$\forall$E} &
\irule{a \N \sep Q\,a}{\exists x.\,Q\,x}{$\exists$I} &
\irule{\exists x.\,Q\,x \sep \dis{a \N,\; Q\,a}{R}}{R}{$\exists$E}\\[2.8ex]
\multicolumn{3}{@{}l}{\textbf{Conditionals}}\\[0.9ex]
\multicolumn{1}{@{}c}{\irule{c \sep a \N}{(\ite{c}{a}{b}) = a}{condI1}} &
\multicolumn{2}{c}{\irule{\neg c \sep b \N}{(\ite{c}{a}{b}) = b}{condI2}}\\[2.4ex]
\multicolumn{1}{@{}c}{\irule{c \sep (\ite{c}{a}{b}) \N}{a \N}{condE1}} &
\multicolumn{2}{c}{\irule{\neg c \sep (\ite{c}{a}{b}) \N}{b \N}{condE2}}\\[2.4ex]
\multicolumn{1}{@{}c}{\irule{(\ite{c}{a}{b}) \N}{c \B}{condE3}} &
\multicolumn{2}{c}{\irule{c \sep Q\,(\ite{c}{a}{b})}{Q\,a}{cond-then-E}}\\[2.4ex]
\multicolumn{1}{@{}c}{\irule{c \sep Q\,a}{Q\,(\ite{c}{a}{b})}{cond-then-I}} &
\multicolumn{2}{c}{\irule{\neg c \sep Q\,b}{Q\,(\ite{c}{a}{b})}{cond-else-I}}\\[2.8ex]
\multicolumn{3}{@{}l}{\textbf{Definitions}}\\[0.9ex]
\irule{a := b \sep Q\,b}{Q\,a}{defE} &
\irule{a := b \sep Q\,a}{Q\,b}{defI} &
\end{tabular}
\caption{Inference rules of \QGA.  Each conditional rule has a formula-level
twin with $=$ replaced by $\longleftrightarrow$ and $\mathsf{N}$ by
$\mathsf{B}$, and each ``then'' rule has an ``else'' twin.}
\label{tab:qga-rules}
\end{table}

\section{\BGA{} arithmetized in \QGA}
\label{sec:bga}

\texttt{BGA\_on\_GA.thy} encodes basic grounded arithmetic as \QGA{} natural
numbers and proves its proof checker sound with respect to the encoded
semantics.  The file is a tower of locales, so the whole argument is carried
out under a fixed, small set of assumptions about the encoding and no concrete
G\"odel numbering appears in it.

This section walks through the locales in order and explains every constant
they fix.  Throughout, $\Gamma \seq \varphi$ is a \BGA{} judgment, $\pi$ and
$\rho$ range over encoded proofs, $A$ over assignments, and $D$ over the
background definition list.  Notation like $q[v_i \mapsto a]$ abbreviates
\isa{subst\_F q i a}.

\subsection{The system being encoded}

\BGA{} has terms and formulas, and no logical connectives at all.

\[
\begin{array}{lcl}
t &::=& v_j \mid 0 \mid \Suc{t} \mid \Pre{t} \mid t\,0?\,t : t \mid d_i(t, t)\\
\varphi &::=& t = t \mid t \neq t
\end{array}
\]

\noindent
A variable $v_j$ is a natural number $j$, and its value under an assignment is
the $j$th entry of the assignment list.  \BGA{} has no binders, so nothing
scopes over a variable and the index is simply a position.  Conditional
evaluation is if-zero-else, which keeps terms and formulas from being mutually
recursive.  Function application is restricted to two arguments, and the
functions come from a fixed background list $D$, where position $i$ holds the
body of $d_i(v_0, v_1)$.  Disequality is a primitive, so $a \neq b$ is not the
negation of $a = b$; \BGA{} has no negation to form.  A judgment is a sequent
$\Gamma \seq \varphi$ with $\Gamma$ a finite list of formulas.

\subsection{Codes}

There is one \QGA{} type, \isa{num}.  The type synonyms below are
documentation.

\thy{bga-types}

\noindent
A judgment is the pair \isa{⟨G, c⟩}, written \isa{G ⊩ c}, with \isa{hyp\_of}
and \isa{conc\_of} as projections.  A hypothesis set is a list of formulas and
a proof is a list of judgments, read head first, each justified by the
judgments after it.

\subsection{The locale tower}
\label{sec:tower}

\begin{center}
\begin{tabular}{@{}lll@{}}
\toprule
\textbf{Locale} & \textbf{Extends} & \textbf{Fixes}\\
\midrule
\isa{suff\_syntax}    & ---                   & \isa{mk\_eq}, \isa{mk\_neq}, \isa{dfns}, \isa{is\_valid\_proof}\\
\isa{suff\_semantics} & \isa{suff\_syntax}    & \isa{evals}, \isa{sat\_fm}, \isa{sat\_hyp}\\
\isa{consistent}      & \isa{suff\_semantics} & ---\\
\midrule
\isa{bga\_encoding}   & ---                   & the six encoder operations, the three\\
                      &                       & freshness predicates\\
\isa{bga\_dfns}       & \isa{bga\_encoding}   & \isa{dfns}, \isa{dfn\_is}\\
\isa{bga\_fuel\_semantics} & \isa{bga\_dfns}  & \isa{eval\_fuel}\\
\isa{bga\_subst}      & \isa{bga\_encoding}   & \isa{subst\_T}, \isa{subst\_F}, \isa{subst\_body}\\
\isa{bga\_fuel\_subst\_semantics} & \isa{bga\_fuel\_semantics}, & \isa{asn\_put}\\
                      & \isa{bga\_subst}      & \\
\midrule
\isa{bga\_subst\_rule}  & \isa{bga\_subst}    & \isa{check\_template}, \isa{rep\_vars\_T},\\
                        &                     & \isa{rep\_vars\_F}, \isa{find\_phi}, \isa{find\_eq},\\
                        &                     & \isa{check\_subst}\\
\isa{bga\_eq\_rule}     & \isa{bga\_encoding} & \isa{check\_eq\_rules}, \isa{check\_neq\_rules}\\
\isa{bga\_ind\_rule}    & \isa{bga\_subst\_rule} & \isa{check\_ind\_template},\\
                        &                     & \isa{find\_ind\_base}, \isa{check\_ind}\\
\isa{bga\_struct\_rule} & \isa{bga\_encoding} & \isa{find\_cut}, \isa{check\_cut},\\
                        &                     & \isa{find\_struct}, \isa{check\_struct}\\
\isa{bga\_app\_rule}    & \isa{bga\_subst\_rule}, & \isa{app\_try}, \isa{app\_y}, \isa{app\_x},\\
                        & \isa{bga\_dfns}     & \isa{app\_d}, \isa{check\_app}\\
\isa{bga\_proof\_check} & the five rule locales & \isa{valid\_step}, \isa{check\_list},\\
                        &                     & \isa{is\_valid\_proof}\\
\midrule
\isa{bga\_fuller}       & \isa{bga\_fuel\_subst\_semantics}, & ---\\
                        & \isa{bga\_proof\_check} & \\
\bottomrule
\end{tabular}
\end{center}

\noindent
Two locales are defined in the file and not used by \isa{bga\_fuller}:
\isa{bga\_semantics}, which fixes the exact evaluator \isa{eval}, and
\isa{bga\_subst\_semantics}, which combines it with substitution.  They are
described in \S\ref{sec:fuel} because the exact evaluator is the clearest
statement of what the encoded semantics means.

Almost every constant is fixed with an assumption of the form \isa{f x := body}
rather than being defined.  A locale cannot make a recursive definition, and in
\QGA{} it does not need to: fixing a constant together with its \isa{:=}
equation is exactly what \S\ref{sec:defmech} admits, and the equation is all
the proofs ever use.

\subsection{The consistency skeleton}
\label{sec:skeleton}

Before any encoding is fixed, the shape of the final argument is stated.

\thy{bga-suff-syntax}

\begin{center}
\begin{tabular}{@{}ll@{}}
\toprule
\isa{mk\_eq :: tm ⇒ tm ⇒ fm} & builds the code of $a = b$ from the codes of $a$ and $b$\\
\isa{mk\_neq :: tm ⇒ tm ⇒ fm} & the same for $a \neq b$\\
\isa{dfns :: dfn} & the code of the background definition list $D$\\
\isa{is\_valid\_proof :: pf ⇒ jdg ⇒ o} & ``$\pi$ is a proof of $J$''\\
\bottomrule
\end{tabular}
\end{center}

\noindent
The assumptions are only habeas quid: encoded formulas are numbers, and
\isa{is\_valid\_proof p J} is a grounded formula whenever \isa{p} and \isa{J}
are numbers.  The second is the grounded-logic form of decidability of proof
checking, and it is what makes the concluding proof by contradiction available.

\thy{bga-suff-semantics}

\begin{center}
\begin{tabular}{@{}ll@{}}
\toprule
\isa{evals :: tm ⇒ asn ⇒ val ⇒ o} & ``$t$ evaluates to $r$ under $A$''\\
\isa{sat\_fm :: fm ⇒ asn ⇒ o} & ``the formula coded by $f$ holds under $A$''\\
\isa{sat\_hyp :: hyp ⇒ asn ⇒ o} & ``every formula in $\Gamma$ holds under $A$''\\
\bottomrule
\end{tabular}
\end{center}

\noindent
Three assumptions: the empty hypothesis list is satisfied by every assignment,
evaluation is deterministic, and a satisfied encoded equation yields a common
value for its two sides while a satisfied encoded disequality yields two
distinct values.

\thy{bga-consistent}

\noindent
\isa{consistent} adds soundness of the checker and derives syntactic
consistency.  The derivation is short.  Assume both proofs exist.  The
statement is grounded by \isa{proof\_bool}, so \isa{contradiction} applies.
Soundness under the empty hypothesis list and the assignment \isa{0} gives
\isa{sat\_fm (mk\_eq a b) 0} and \isa{sat\_fm (mk\_neq a b) 0}.  The first
produces a value \isa{q} reached by both \isa{a} and \isa{b}, the second
produces \isa{x} and \isa{y} with \isa{x ≠ y}, determinism forces
\isa{x = q = y}, and \isa{exF} closes the goal.  The habeas quid premises
\isa{a N}, \isa{b N}, \isa{p1 N} and \isa{p2 N} are needed before the formula
being contradicted can be shown grounded.

\subsection{Encoded syntax}

Terms and formulas are tagged pairs, with six term tags and two formula tags.

\thy{bga-tags}

\begin{center}
\begin{tabular}{@{}llll@{}}
\toprule
Tag & Constructor & Notation & Payload\\
\midrule
\isa{T\_VAR}  & variable      & $v_j$              & $j$\\
\isa{T\_ZERO} & zero          & $0$                & ignored\\
\isa{T\_SUC}  & successor     & $\Suc{a}$          & $a$\\
\isa{T\_PRED} & predecessor   & $\Pre{a}$          & $a$\\
\isa{T\_IFZ}  & if-zero-else  & $c\;0?\,a:b$       & $\langle c, \langle a, b\rangle\rangle$\\
\isa{T\_APP}  & application   & $d_i(a,b)$         & $\langle i, \langle a, b\rangle\rangle$\\
\midrule
\isa{F\_EQ}   & equality      & $a = b$            & $\langle a, b\rangle$\\
\isa{F\_NEQ}  & disequality   & $a \neq b$         & $\langle a, b\rangle$\\
\bottomrule
\end{tabular}
\end{center}

\subsection{\isa{bga\_encoding}}
\label{sec:encoder-interface}

\thy{bga-enc-fixes}

\begin{center}
\begin{tabular}{@{}ll@{}}
\toprule
\isa{tag\_T :: tm ⇒ tmtag} & the constructor tag of a term code\\
\isa{load\_T :: tm ⇒ tm}   & the payload of a term code\\
\isa{pack\_T :: tmtag ⇒ tm ⇒ tm} & builds a term code from tag and payload\\
\isa{tag\_F}, \isa{load\_F}, \isa{pack\_F} & the same three for formula codes\\
\isa{fresh\_T :: num ⇒ tm ⇒ o} & every variable in the term has index below $k$\\
\isa{fresh\_F :: num ⇒ fm ⇒ o} & the same for a formula\\
\isa{fresh\_H :: num ⇒ hyp ⇒ o} & the same for every formula in a list\\
\bottomrule
\end{tabular}
\end{center}

\noindent
Freshness is defined by recursion over the tag, and in mathematical notation it
reads
\[
\fn{fresh}_T(k, t) =
\begin{cases}
j < k                                        & t = v_j\\
\top                                         & t = 0\\
\fn{fresh}_T(k, a)                           & t = \Suc{a} \text{ or } t = \Pre{a}\\
\fn{fresh}_T(k, c) \wedge \fn{fresh}_T(k, a) \wedge \fn{fresh}_T(k, b)
                                             & t = (c\,0?\,a:b)\\
\fn{fresh}_T(k, a) \wedge \fn{fresh}_T(k, b) & t = d_i(a, b)\\
\bot                                         & \text{tag out of range}
\end{cases}
\]
with $\fn{fresh}_F(k, \varphi)$ the conjunction over the two sides of
$\varphi$ and $\fn{fresh}_H(k, \Gamma)$ the conjunction over the members of
$\Gamma$.  The last line is the treatment of a number whose tag is not one of
the six: it is fresh for no $k$.  Freshness is used once, in the induction rule
of \S\ref{sec:ind-rule}, to guarantee that the induction variable does not
occur in the hypotheses.

The assumptions on the interface are the following, and they are the only
properties of the encoding used anywhere in the development.

\thy{bga-enc-axioms}

\begin{center}
\begin{tabular}{@{}ll@{}}
\toprule
habeas quid & \isa{pack}, \isa{tag} and \isa{load} preserve \isa{N}\\
injectivity & \isa{tag (pack t L) = t} and \isa{load (pack t L) = L}\\
retraction & \isa{pack (tag f) (load f) = f}\\
decrease & \isa{f ≠ 0 ⟹ load f < f = 1}, payloads are strictly smaller\\
base case & \isa{tag\_T 0 = T\_ZERO}, the code $0$ is the term $0$\\
monotonicity & \isa{pack} is monotone in its payload\\
\bottomrule
\end{tabular}
\end{center}

\begin{remarkx}[Injective, not surjective]
What the concrete encoding of \S\ref{sec:encode} provides, and all that is
used, is an \emph{injective} map from \BGA{} syntax into $\mathbb{N}$.  It is
not surjective: a number whose tag lies outside the constructor range encodes
no \BGA{} term, and no attempt is made to exclude such numbers or to recognize
the well-formed ones.  The assumption that does hold for every number is the
retraction property \isa{pack (tag f) (load f) = f}, which follows from the
Cantor pairing being a bijection on $\mathbb{N}$ and says only that a code can
be taken apart and put back together.  Ill-formed codes are harmless: the
evaluator gives them the application semantics, freshness rejects them, and
soundness never needs them to denote anything.
\end{remarkx}

\noindent
The decrease and monotonicity assumptions are what allow
\isa{strong\_induction} to serve as structural induction over encoded terms.
Decrease gives the descent: the payload of a nonzero code is a strictly smaller
number, so recursion into subterms is bounded numerical recursion.
Monotonicity is used differently, to justify the bounds of the searches in
\S\ref{sec:subst-rule}.  The first lemma proved from them is
\isa{fresh\_T\_bool}, that freshness is a grounded formula.

\subsection{\isa{bga\_dfns}}

\thy{bga-dfns}

\begin{center}
\begin{tabular}{@{}ll@{}}
\toprule
\isa{dfns :: dfn} & the code of the definition list $D$, a list of term codes\\
\isa{dfn\_is :: dfn ⇒ num ⇒ tm ⇒ o} & ``$d$ is in range, $D_d$ is $b$, and $b$ has arity $\leq k$''\\
\bottomrule
\end{tabular}
\end{center}

\noindent
Written out, $\fn{dfn\_is}(d, k, b)$ holds when $d < |D|$, $D_d = b$ and
$\fn{fresh}_T(k, b)$.  The range test is written as a nested conditional rather
than a conjunction so that $\fn{nth}\,d\,D$ is never mentioned when $d$ is out
of range.  The application rule instantiates it at $k = 2$, which is the sense
in which \BGA{} definitions are binary: the body may mention $v_0$ and $v_1$
and nothing else.

\subsection{Semantics}
\label{sec:fuel}

Two evaluators appear in the file.  The direct one, fixed by
\isa{bga\_semantics}, is the natural reading of the \BGA{} big-step semantics,
and writing it down at all depends on \QGA{} tolerating divergent definitions.

\thy{bga-eval}

\noindent
In mathematical notation, with $A$ an assignment,
\[
\fn{eval}(t, A) =
\begin{cases}
\fn{nth}\,j\,A                     & t = v_j\\
0                                  & t = 0\\
\Suc{\fn{eval}(a, A)}              & t = \Suc{a}\\
\Pre{\fn{eval}(a, A)}              & t = \Pre{a}\\
\fn{eval}(a, A) \text{ or } \fn{eval}(b, A)
                                   & t = (c\,0?\,a:b), \text{ on } \fn{eval}(c, A) = 0\\
\fn{eval}(D_i,\; [\fn{eval}(a, A),\, \fn{eval}(b, A)])
                                   & t = d_i(a, b)
\end{cases}
\]
The application case in the sources is wrapped in two guards of the form
\isa{eval a A = eval a A}.  Those are habeas quid tests: the arguments are
evaluated first, and the body is entered only if both terminate, which makes
the encoded semantics call by value.

The consistency argument runs through the step-indexed variant instead, fixed
by \isa{bga\_fuel\_semantics}.  Its result is offset by one, so that
$\fn{eval\_fuel}(k, t, A) = 0$ means ``no result within budget $k$'' and
$\fn{eval\_fuel}(k, t, A) = 1 + v$ means ``the value is $v$''.  The offset is
what lets a single \isa{num} carry both outcomes, at the cost of an explicit
$\Pre{\cdot}$ on every recursive result.

\thy{bga-eval-fuel}
\thy{bga-evals}

\noindent
\isa{evals t A r} says that $t$ evaluates to $r$ under $A$ for some budget:
\[
\fn{evals}(t, A, r) \;\;\equiv\;\; \exists k.\; \fn{eval\_fuel}(k, t, A) = \Suc{r}.
\]

Satisfaction comes in the same two forms.  For the direct evaluator,
\isa{bga\_semantics} defines

\thy{bga-sat}

\noindent
that is, $\fn{sat}(\varphi, A)$ compares $\fn{eval}$ of the two sides, with the
comparison chosen by the formula tag, and $\fn{sat\_hyp}$ quantifies over the
members of the list.  For the step-indexed evaluator, satisfaction of an
equation also has to require both sides to converge:

\thy{bga-sat-fuel}
\thy{bga-sat-hyp-fuel}

\[
\fn{sat\_fuel}(\varphi, A) \;\;\equiv\;\;
\exists x, y.\;\; \fn{evals}(a, A, x) \wedge \fn{evals}(b, A, y) \wedge
\begin{cases} x = y & \text{tag } \varphi = \fn{F\_EQ}\\
              x \neq y & \text{otherwise}\end{cases}
\]
where $a$ and $b$ are the two sides of $\varphi$.  A formula whose sides
diverge is unsatisfied under every assignment, which is the property the
consistency argument ultimately rests on.

\subsection{\isa{bga\_subst} and \isa{asn\_put}}

\thy{bga-subst}
\thy{bga-asn-put}

\begin{center}
\begin{tabular}{@{}ll@{}}
\toprule
\isa{subst\_T :: tm ⇒ tm ⇒ tm ⇒ tm} & $t[v_j \mapsto v]$, replace variable $j$ inside a term\\
\isa{subst\_F :: fm ⇒ tm ⇒ tm ⇒ fm} & the same inside a formula, side by side\\
\isa{subst\_body :: tm ⇒ tm ⇒ tm ⇒ tm} & $b[v_0 \mapsto x,\, v_1 \mapsto y]$, simultaneous\\
\isa{asn\_put :: asn ⇒ num ⇒ val ⇒ asn} & $A[i \mapsto v]$, update an assignment\\
\bottomrule
\end{tabular}
\end{center}

\noindent
All three substitutions recurse over the tag and rebuild with \isa{pack\_T}.
Two details are worth noting.  In the \isa{T\_ZERO} case \isa{subst\_T} returns
\isa{pack\_T T\_ZERO 0} rather than the original code, which normalizes the
ignored payload of a zero.  In the \isa{T\_APP} case the definition index
\isa{cpx (load\_T t)} is carried through untouched, since it is not a term.
\isa{subst\_body} is the two-variable version used by the application rule, and
it exists separately because that rule substitutes both arguments at once.
\isa{asn\_put A i v} walks down $i$ entries and replaces the one it finds,
padding with zeros when $i$ runs past the end of the list, so that an
assignment is a total map from indices to values.

\subsection{The proof rules of \BGA}
\label{sec:rules}

The rules are given below in sequent form, and each is realized as one branch
of the encoded checker.  Here $\Gamma$ is a list of encoded formulas,
$\varphi$ an encoded formula, and $\varphi[a]$ a formula with occurrences of
$a$ marked, which is realized by the template search of
\S\ref{sec:subst-rule}.  The turnstile is the \BGA{} sequent arrow, written
\isa{⊩} in the sources.

\begin{rules}
\iaxiom{\Gamma, \varphi \seq \varphi}{hyp} &
\iaxiom{\Gamma \seq 0 = 0}{0I}\\[2.4ex]
\irule{\Gamma \seq a = b}{\Gamma \seq b = a}{=S} &
\irule{\Gamma \seq a = b \sep \Gamma \seq \varphi[a]}{\Gamma \seq \varphi[b]}{=E}\\[2.4ex]
\irule{\Gamma \seq a = b}{\Gamma \seq \Suc{a} = \Suc{b}}{S=I} &
\irule{\Gamma \seq \Suc{a} = \Suc{b}}{\Gamma \seq a = b}{S=E}\\[2.4ex]
\irule{\Gamma \seq a \N}{\Gamma \seq \Pre{\Suc{a}} = a}{PI} &
\irule{\Gamma \seq a \N}{\Gamma \seq \Suc{a} \neq 0}{S{\neq}0I}\\[2.4ex]
\irule{\Gamma \seq c = 0 \sep \Gamma \seq a \N}{\Gamma \seq (c\,0?\,a:b) = a}{?I1} &
\irule{\Gamma \seq c \neq 0 \sep \Gamma \seq b \N}{\Gamma \seq (c\,0?\,a:b) = b}{?I2}\\[2.4ex]
\irule{\Gamma \seq a \neq b}{\Gamma \seq b \neq a}{{\neq}S} &
\irule{\Gamma \seq a \neq b}{\Gamma \seq \Suc{a} \neq \Suc{b}}{{\neq}SI}\\[2.4ex]
\irule{\Gamma \seq \Suc{a} \neq \Suc{b}}{\Gamma \seq a \neq b}{S{\neq}E} &
\irule{\Gamma \seq a \sep \Gamma, a \seq c}{\Gamma \seq c}{cut}\\[2.4ex]
\irule{\Gamma' \seq \varphi \sep \Gamma' \subseteq \Gamma}{\Gamma \seq \varphi}{wk} &
\irule{\begin{array}{c}
        \Gamma \seq a \N \sep \Gamma \seq \varphi[0]\\
        \Gamma, v_i \N, \varphi[v_i] \seq \varphi[\Suc{v_i}]
       \end{array}}{\Gamma \seq \varphi[a]}{Ind}\\[4.4ex]
\multicolumn{2}{c}{%
\irule{\Gamma \seq a_0 \N \sep \Gamma \seq a_1 \N \sep
       \Gamma \seq \varphi[\,D_d\langle a_0, a_1\rangle\,]}
      {\Gamma \seq \varphi[\,d(a_0, a_1)\,]}{app}}
\end{rules}

\noindent
Termination is not a separate judgment form in \BGA.  $a \N$ abbreviates the
formula $a = a$.  Almost every rule is an introduction rule, and the only
elimination is $=$E.

The checker never applies a rule forwards.  It is given a judgment $J$ and the
tail $\rho$ of the proof, and asks whether $J$ could have been produced by some
rule from judgments occurring in $\rho$.  Every rule with a premise therefore
becomes a search through $\rho$, and the rules whose premises mention a
$\varphi[\cdot]$ become a search through $\rho$ combined with a search for a
template.

\subsection{\isa{bga\_subst\_rule}}
\label{sec:subst-rule}

\begin{center}
\begin{tabular}{@{}ll@{}}
\toprule
\isa{rep\_vars\_T :: tm ⇒ tm ⇒ tm} & replaces every leaf of a term by a given variable\\
\isa{rep\_vars\_F :: fm ⇒ tm ⇒ fm} & the same on both sides of a formula\\
\isa{check\_template} & searches for a template below a bound\\
\isa{find\_phi} & searches $\rho$ for the premise $\Gamma \seq \varphi[a]$\\
\isa{find\_eq} & searches $\rho$ for the equation $\Gamma \seq a = b$\\
\isa{check\_subst :: jdg ⇒ pf ⇒ o} & the rule $=$E as a whole\\
\bottomrule
\end{tabular}
\end{center}

\noindent
The substitution rule is the one place where the checker has to invent
something.  To justify $\Gamma \seq \varphi[b]$ it must produce three things: a
premise $\Gamma \seq \varphi[a]$ somewhere in $\rho$, an equation
$\Gamma \seq a = b$ somewhere in $\rho$, and a \emph{template} $p$, a formula
with a designated variable $v_i$, such that
\[
p[v_i \mapsto a] = \varphi[a] \qquad\text{and}\qquad p[v_i \mapsto b] = \varphi[b].
\]
The template records which occurrences of $a$ are being rewritten.  Nothing in
the conclusion says what it is, so it is searched for.

\paragraph{Choosing the variable.}
The designated variable must not already occur in the judgment, or substituting
into the template would disturb something else.  The index chosen is $J + 1$,
where $J$ is the code of the whole judgment.  This is fresh because a variable
$v_j$ is encoded as $\fn{pack}_T(\fn{T\_VAR}, j)$, which is a pair with second
component $j$, and pairing and packing are monotone in that component, so
$j \leq J$ for every variable occurring anywhere in $J$.  Monotonicity of
\isa{pack} in \S\ref{sec:encoder-interface} is assumed for this argument.

\paragraph{Bounding the search.}
Templates are enumerated downwards from
$\fn{rep\_vars}_F(\varphi[b], J+1)$, the formula obtained by replacing every
leaf of the conclusion by $v_{J+1}$:
\[
\fn{rep\_vars}_T(t, i) =
\begin{cases}
v_i                                        & t = v_j \text{ or } t = 0\\
\Suc{\fn{rep\_vars}_T(a, i)}               & t = \Suc{a}\\
\Pre{\fn{rep\_vars}_T(a, i)}               & t = \Pre{a}\\
(\fn{rep\_vars}_T(c, i)\,0?\,\fn{rep\_vars}_T(a, i) : \fn{rep\_vars}_T(b, i))
                                           & t = (c\,0?\,a:b)\\
d_j(\fn{rep\_vars}_T(a, i), \fn{rep\_vars}_T(b, i))
                                           & t = d_j(a, b)
\end{cases}
\]
This is the template that abstracts \emph{every} leaf at once.  Its code
bounds the code of every other candidate with the same constructor skeleton,
since such a candidate differs from it only by leaving some leaves
unabstracted, and a leaf is numerically smaller than $v_{J+1}$.
\isa{check\_template} then counts $p$ down from that bound to zero and tests
each value:
\[
\fn{check\_template}(f, \varphi, a, b, p, i) \;\;\Longleftrightarrow\;\;
\exists q \leq p.\;\; q[v_i \mapsto a] = \varphi \;\wedge\; q[v_i \mapsto b] = f.
\]

\thy{bga-check-template}
\thy{bga-rep-vars-t}

\begin{remarkx}[The bound matters for completeness only]
Soundness of the checker does not depend on the bound being correct.
\isa{check\_template\_sound\_fuel} assumes that the search reported success and
extracts the witness $q$ it found; a bound that was too small would make the
checker reject proofs it ought to accept, which is a completeness failure, and
would leave soundness untouched.  Nothing in this development proves
completeness.
\end{remarkx}

\paragraph{Searching the proof.}
With the template search in place, the two remaining searches walk the proof
list.  In mathematical notation, writing $K \in \rho$ for membership in the
proof list,
\begin{align*}
\fn{find\_phi}(J, a, b, \pi) &\;\Longleftrightarrow\;
  \exists K \in \pi.\;\; \fn{hyp}(K) = \fn{hyp}(J) \;\wedge\;
  \fn{check\_template}(\fn{conc}(J), \fn{conc}(K), a, b, \beta_J, J+1),\\
\fn{find\_eq}(J, \rho, \pi) &\;\Longleftrightarrow\;
  \exists K \in \pi.\;\; \fn{hyp}(K) = \fn{hyp}(J) \;\wedge\;
  \fn{conc}(K) \text{ is } (a = b) \;\wedge\; \fn{find\_phi}(J, a, b, \rho),\\
\fn{check\_subst}(J, \rho) &\;\Longleftrightarrow\; \fn{find\_eq}(J, \rho, \rho),
\end{align*}
where $\beta_J = \fn{rep\_vars}_F(\fn{conc}(J), J+1)$ is the template bound.
Both searches carry two proof arguments.  The one named \isa{ptr} is the
part still to be scanned and shrinks on each step, and the one named
\isa{rest} is the whole tail of the proof and stays fixed, because the
premise being looked for may sit anywhere in it.  The hypothesis test
$\fn{hyp}(K) = \fn{hyp}(J)$ is what keeps the rule from silently changing the
context.

\thy{bga-find-phi}
\thy{bga-find-eq}

\subsection{\isa{bga\_eq\_rule}}

\begin{center}
\begin{tabular}{@{}ll@{}}
\toprule
\isa{check\_eq\_rules :: hyp ⇒ tm ⇒ tm ⇒ tmtag ⇒ tmtag ⇒ pf ⇒ o}
  & the seven rules concluding an equation\\
\isa{check\_neq\_rules :: hyp ⇒ tm ⇒ tm ⇒ tmtag ⇒ tmtag ⇒ pf ⇒ o}
  & the four rules concluding a disequation\\
\bottomrule
\end{tabular}
\end{center}

\noindent
Both take the two sides and their tags as separate arguments, already projected
out of the conclusion by \isa{valid\_step}, so that the branches can test tags
without recomputing them.

\thy{bga-eq-rules}

\noindent
The seven branches are, in order, \textsc{0I}, \textsc{=S}, \textsc{S=I},
\textsc{S=E}, \textsc{PI}, \textsc{?I2} and \textsc{?I1}.  Each is a membership
test on \isa{rest}, so a judgment is justified only by judgments occurring
later in the list, which is what makes \isa{check\_list} a well-founded
recursion.  Reading the fifth branch as an example: the conclusion is
$\Gamma \seq \Pre{u} = b$, the branch checks that $u$ is $\Suc{b}$ by
inspecting tags and payloads, and it requires $\Gamma \seq b = b$ to appear in
the proof, which is the habeas quid premise $b \N$ of \textsc{PI}.

\thy{bga-neq-rules}

\noindent
and dually \textsc{$\neq$S}, \textsc{S$\neq$0I}, \textsc{$\neq$SI} and
\textsc{S$\neq$E}.

\subsection{\isa{bga\_ind\_rule}}
\label{sec:ind-rule}

\begin{center}
\begin{tabular}{@{}ll@{}}
\toprule
\isa{check\_ind\_template} & tests one candidate template for the induction rule\\
\isa{find\_ind\_base} & searches $\rho$ for the base case\\
\isa{check\_ind :: jdg ⇒ pf ⇒ o} & the rule \textsc{Ind} as a whole\\
\bottomrule
\end{tabular}
\end{center}

\noindent
Induction needs the same kind of template, and three premises rather than one.
To justify $\Gamma \seq \varphi[a]$ it looks for a template $q$ with
\[
q[v_i \mapsto a] = \varphi[a], \qquad
q[v_i \mapsto 0] \text{ present in } \rho, \qquad
(\,v_i = v_i,\; q,\; \Gamma \seq q[v_i \mapsto \Suc{v_i}]\,) \text{ present in } \rho,
\]
that is,
\begin{align*}
\fn{check\_ind\_template}(f, \varphi, a, \Gamma, p, i, \rho)
  &\;\Longleftrightarrow\; \exists q \leq p.\;\;
   q[v_i \mapsto a] = f \;\wedge\; q[v_i \mapsto 0] = \varphi\\
  &\qquad\wedge\;\;
   \bigl(v_i = v_i,\, q,\, \Gamma \seq q[v_i \mapsto \Suc{v_i}]\bigr) \in \rho.
\end{align*}

\thy{bga-ind-template}
\thy{bga-check-ind}

\noindent
\isa{check\_ind} guards the search with two conditions.  The habeas quid
premise $\Gamma \seq a \N$, encoded as the judgment $\Gamma \seq a = a$, must
be in the proof, and $\fn{fresh}_H(J+1, \Gamma)$ must hold, so that the
variable used to build the template does not occur in the hypotheses.  The
step case sought in the proof carries both the induction hypothesis $q$ and the
termination assumption $v_i = v_i$ on the induction variable, which is what the
semantic induction lemma of \S\ref{sec:proof} consumes.  \isa{find\_ind\_base}
scans $\rho$ for the base case in the same style as \isa{find\_phi}.

\subsection{\isa{bga\_struct\_rule}}

\begin{center}
\begin{tabular}{@{}ll@{}}
\toprule
\isa{find\_cut} & searches $\rho$ for the cut formula\\
\isa{check\_cut :: jdg ⇒ pf ⇒ o} & the rule \textsc{cut}\\
\isa{find\_struct} & searches $\rho$ for the weakened judgment\\
\isa{check\_struct :: jdg ⇒ pf ⇒ o} & the rule \textsc{wk}\\
\bottomrule
\end{tabular}
\end{center}

\thy{bga-struct}

\begin{align*}
\fn{check\_cut}(J, \rho) &\;\Longleftrightarrow\;
  \exists K \in \rho.\;\; \fn{hyp}(K) = \fn{hyp}(J) \;\wedge\;
  \bigl(\fn{conc}(K),\, \fn{hyp}(J) \seq \fn{conc}(J)\bigr) \in \rho,\\
\fn{check\_struct}(J, \rho) &\;\Longleftrightarrow\;
  \exists K \in \rho.\;\; \fn{conc}(K) = \fn{conc}(J) \;\wedge\;
  \fn{hyp}(K) \subseteq \fn{hyp}(J).
\end{align*}

\noindent
The hypothesis rule is not a separate function.  It is the first branch of
\isa{valid\_step}, the test \isa{mem (conc\_of J) (hyp\_of J)}.

\subsection{\isa{bga\_app\_rule}}

\begin{center}
\begin{tabular}{@{}ll@{}}
\toprule
\isa{app\_try} & tests one triple $(d, x, y)$\\
\isa{app\_y} & counts $y$ down\\
\isa{app\_x} & counts $x$ down\\
\isa{app\_d} & counts the definition index $d$ down\\
\isa{check\_app :: jdg ⇒ pf ⇒ o} & the rule \textsc{app} as a whole\\
\bottomrule
\end{tabular}
\end{center}

\thy{bga-app}

\noindent
Application unfolds a definition.  To justify $\Gamma \seq \varphi[d(x,y)]$ the
checker needs $\Gamma \seq x \N$, $\Gamma \seq y \N$ and a premise
$\Gamma \seq \varphi[\,D_d[v_0 \mapsto x, v_1 \mapsto y]\,]$:
\begin{align*}
\fn{app\_try}(J, d, x, y, \rho) \;\Longleftrightarrow\;\;
  & \fn{dfn\_is}(d, 2, D_d)
    \;\wedge\; (\fn{hyp}(J) \seq x = x) \in \rho
    \;\wedge\; (\fn{hyp}(J) \seq y = y) \in \rho\\
  &\wedge\; \fn{find\_phi}\bigl(J,\; D_d[v_0 \mapsto x, v_1 \mapsto y],\; d(x,y),\; \rho\bigr).
\end{align*}
Note that \isa{find\_phi} is reused here with the two sides of the rewrite
being the unfolded body and the application, so unfolding a definition is
treated as one more instance of rewriting inside a template.

None of $d$, $x$ and $y$ is determined by the conclusion, so all three are
enumerated:
\[
\fn{check\_app}(J, \rho) \;\Longleftrightarrow\;
\exists d < |D|.\;\; \exists x \leq \fn{conc}(J).\;\; \exists y \leq \fn{conc}(J).\;\;
\fn{app\_try}(J, d, x, y, \rho).
\]
The three nested countdowns \isa{app\_d}, \isa{app\_x} and \isa{app\_y} are
that triple search written as recursion, since \QGA{} has no bounded
quantifier.  The bound $\fn{conc}(J)$ on the arguments is the same kind of
bound as in \S\ref{sec:subst-rule}: the code of a subterm is at most the code
of the formula containing it, and as before it is completeness rather than
soundness that depends on it.

\subsection{\isa{bga\_proof\_check}}

\begin{center}
\begin{tabular}{@{}ll@{}}
\toprule
\isa{valid\_step :: jdg ⇒ pf ⇒ o} & ``$J$ follows by one rule from judgments in $\rho$''\\
\isa{check\_list :: pf ⇒ o} & every entry follows from the entries after it\\
\isa{is\_valid\_proof :: pf ⇒ jdg ⇒ o} & ``$\pi$ is a valid proof whose head is $J$''\\
\bottomrule
\end{tabular}
\end{center}

\thy{bga-valid-step}

\noindent
\isa{valid\_step} is the disjunction of the branches, tried in order:
membership for \textsc{hyp}, then cut, substitution, induction, application and
weakening, and finally the equality or disequality rules according to the tag
of the conclusion.  Written as a formula,
\begin{align*}
\fn{valid\_step}(J, \rho) \;\Longleftrightarrow\;\;
& \fn{conc}(J) \in \fn{hyp}(J)
  \;\vee\; \fn{check\_cut}(J, \rho)
  \;\vee\; \fn{check\_subst}(J, \rho)\\
&\vee\; \fn{check\_ind}(J, \rho)
  \;\vee\; \fn{check\_app}(J, \rho)
  \;\vee\; \fn{check\_struct}(J, \rho)\\
&\vee\; \bigl(\text{tag } \fn{conc}(J) = \fn{F\_EQ}
   \;?\; \fn{check\_eq\_rules}(\dots) : \fn{check\_neq\_rules}(\dots)\bigr).
\end{align*}
The order is not significant for soundness.  It is a sequence of nested
conditionals rather than a disjunction because each branch has to be evaluated
only when the earlier ones fail.

\thy{bga-check-list}
\thy{bga-is-valid-proof}

\noindent
A proof is valid when every entry follows from the entries after it, and it
proves $J$ when additionally its head is $J$:
\[
\fn{check\_list}(\pi) \;\Longleftrightarrow\;
\pi = \fn{Nil} \;\vee\;
\bigl(\fn{valid\_step}(\fn{hd}\,\pi,\, \fn{tl}\,\pi) \wedge \fn{check\_list}(\fn{tl}\,\pi)\bigr),
\]
\[
\fn{is\_valid\_proof}(\pi, J) \;\Longleftrightarrow\;
\pi \neq \fn{Nil} \;\wedge\; \fn{hd}\,\pi = J \;\wedge\; \fn{check\_list}(\pi).
\]
Reading the list head first, with each entry justified by its tail, is what
makes the recursion terminate and what the soundness induction of
\S\ref{sec:list} follows.

\section{The consistency proof}
\label{sec:proof}

This section follows the argument from the evaluator up to the final theorem.
The intent is to make the shape of the proof and its load-bearing lemmas
visible; the detailed proofs are left to a later instalment.

\subsection{Shape of the argument}

Everything reduces to one implication, called the soundness bridge:

\begin{center}
if \isa{p} checks as a proof of \isa{J}, and the hypotheses of \isa{J} are
satisfied by an assignment \isa{A}, then the conclusion of \isa{J} is satisfied
by \isa{A}.
\end{center}

\noindent
Given it, \S\ref{sec:skeleton} closes the argument in half a page.  Proving it
requires four things, and they are the subsections that follow.

\begin{enumerate}
\item The step-indexed evaluator determines at most one value, so that
      \isa{evals} behaves as a partial function and satisfaction is well
      behaved.
\item Every branch of \isa{valid\_step} preserves satisfaction, assuming that
      every judgment in the tail of the proof does.
\item That local statement can be lifted along the proof list by induction.
\item Satisfaction of an encoded equation and of an encoded disequality can be
      unpacked into statements about values.
\end{enumerate}

\noindent
A remark on technique before the details.  \QGA{} induction instantiates a
schematic predicate \isa{Q}, so a statement can only be generalized before
inducting if the generalization is part of the object-level formula.  Several
proofs in this file are therefore stated internally in the form
\isa{∀u. ∀A2. …}, or as nested implications \isa{… ⟶ (… ⟶ …)}, and unpacked
afterwards with \isa{forallE} and \isa{implE}.  \isa{eval\_fuel\_N} is the first
instance: the induction on fuel runs on \isa{∀u. ∀A2. eval\_fuel k u A2 N} and
the desired statement is recovered by two instantiations.

Nesting an implication inside a motive has a price.  $\longrightarrow$ introduction
requires its antecedent to be grounded, so every antecedent of such a motive
needs a groundedness lemma.  The checkers all have one.  Satisfaction does not,
and \S\ref{sec:satat} is what supplies it.

\subsection{The evaluator is a partial function}

The fuel evaluator is total as a function into \isa{num}, which is the
statement \isa{eval\_fuel\_N: ⟦k N; t N; A N⟧ ⟹ eval\_fuel k t A N}, and $0$ is
its ``no result'' value.  Success is preserved when the budget grows.

\thy{bga-fuel-add}

\noindent
This is proved by induction on the increment, from the one-step version
\isa{eval\_fuel\_success\_suc}.  Two successful runs with different budgets
agree, since both can be raised to the common budget \isa{k + l}.

\thy{bga-fuel-unique}

\noindent
The existential form follows, and it is the assumption \isa{evals\_det} that
the abstract locale of \S\ref{sec:skeleton} demanded.

\thy{bga-evals-functional}

\noindent
The remaining facts about the evaluator are one lemma per constructor, in two
directions each: from a successful run one can read off the value of the
subterms (\isa{eval\_fuel\_suc\_value}, \isa{eval\_fuel\_pred\_value},
\isa{eval\_fuel\_ifz\_zero\_value}, \isa{eval\_fuel\_app\_value}), and from
successful runs of the subterms one can build a successful run of the whole
(\isa{eval\_fuel\_sucI} and its relatives).  Together they let \isa{evals} be
used as an inductively defined big-step relation, which is how the rule
soundness proofs consume it.

\subsection{Soundness of a single step}

The central local statement is the following.  It says that a checked step is
sound relative to the tail of the proof, and it is what the induction of
\S\ref{sec:list} will feed itself.

\thy{bga-valid-step-sound}

\noindent
Its proof is a case analysis over the branches of \isa{valid\_step}, using
\isa{cases\_bool} on each guard in turn, with one lemma per branch.  All the
branch lemmas have the same premises as above and the same conclusion; they
differ only in which checker they assume.

\begin{center}
\begin{tabular}{@{}ll@{}}
\toprule
\textbf{Branch} & \textbf{Lemma}\\
\midrule
hypothesis          & inline, from \isa{sat\_hyp\_fuel\_mem}\\
cut                 & \isa{check\_cut\_sound\_fuel}\\
substitution of equals & \isa{check\_subst\_sound\_fuel}\\
induction           & \isa{check\_ind\_sound\_fuel}\\
application         & \isa{check\_app\_sound\_fuel}\\
weakening           & \isa{check\_struct\_sound\_fuel}\\
the seven equality rules   & \isa{check\_eq\_rules\_sound\_fuel}\\
the four disequality rules & \isa{check\_neq\_rules\_sound\_fuel}\\
\bottomrule
\end{tabular}
\end{center}

\noindent
Two of these carry the real content.

\paragraph{The template lemma.}
Substitution of equals, application and induction all reduce to the same
semantic fact: if a template \isa{p} instantiated at \isa{a} is satisfied, and
\isa{a} and \isa{b} evaluate to a common value, then the same template
instantiated at \isa{b} is satisfied.

\thy{bga-template-sound}

\noindent
It rests on a pair of substitution lemmas that relate syntactic substitution to
updating the assignment.

\thy{bga-subst-down}
\thy{bga-subst-up}

\noindent
Given \isa{a} and \isa{b} that both evaluate to \isa{v}, the first pushes
satisfaction of \isa{p[i ↦ a]} under \isa{A} down to satisfaction of \isa{p}
under \isa{asn\_put A i v}, and the second lifts it back up to
\isa{p[i ↦ b]}.  The search functions are then sound because they report
success only when such a template exists, which is the content of
\isa{check\_template\_sound\_fuel} and of

\thy{bga-find-phi-sound}

\noindent
The premise \isa{prev} is the induction hypothesis of \S\ref{sec:list} in the
form the search needs: every judgment occurring in \isa{rest} is sound.  The
premise \isa{sub} records that the pointer walks a sublist of \isa{rest}, which
is what keeps the recursion inside the part of the proof already known to be
sound.  Application is the same lemma applied to the template
$\varphi[\,\mathrm{body}_d[v_0 \mapsto x, v_1 \mapsto y]\,]$, with the two
habeas quid premises supplying the values that \isa{subst\_body} needs.

\paragraph{The induction rule.}
Soundness of the induction rule is the one place where a semantic induction is
performed, over the value of the term the rule inducts on.

\thy{bga-nat-ind-sound}

\noindent
Read semantically: if the hypothesis list is satisfied, the base instance
$p[v_i \mapsto 0]$ is satisfied, and the step is satisfied for every value
\isa{m} put into the assignment at index \isa{i}, then the instance at \isa{a}
is satisfied whenever \isa{a} evaluates to \isa{n}.  The proof establishes
\isa{⋀m. m N ⟹ sat\_fuel p (asn\_put A i m)} by \QGA{} induction on \isa{m},
with the base case supplied by \isa{sat\_fuel\_subst\_FD} applied to the
encoded zero, and then instantiates it at \isa{n}.  The freshness premise
\isa{fresh\_H i G} is what
allows the assignment to be modified at \isa{i} without disturbing satisfaction
of the hypotheses.  \isa{check\_ind\_sound\_fuel} then supplies the three
premises from the corresponding judgments found in the proof list.

\subsection{From steps to proofs}
\label{sec:list}

The lifting step is a dedicated induction principle over proof lists.

\thy{bga-check-list-induct}

\noindent
It is stated for \isa{sat\_fuel} and \isa{sat\_hyp\_fuel}.  Its only caller
instantiates it at those two, and the indexed predicates it uses internally are
defined in terms of them.

Internally it is proved by \isa{list\_induct} on a statement of the form
\begin{center}
\isa{∀A. ∀K. ∀k. check\_list pf ⟶ (mem K pf ⟶ (sat\_hyp\_at k (hyp\_of K) A ⟶ sat\_fuel (conc\_of K) A))}.
\end{center}
The predicate being inducted on has to be a single formula, so the
generalizations over the assignment, over the judgment and over the budget sit
inside it as \QGA{} quantifiers.

The three antecedents are why the budget is explicit.  Each one has to be
grounded before $\longrightarrow$ can be introduced, and \isa{check\_list\_bool},
\isa{mem\_bool} and \isa{sat\_hyp\_at\_bool} supply exactly that at the three
levels.  The conclusion needs no such treatment, since $\longrightarrow$ introduction
constrains only the antecedent, which is why \isa{sat\_fuel} appears there
unindexed.  With \isa{sat\_hyp\_fuel} in the antecedent instead, the introduction
would be blocked.

The indexed predicate never leaves the proof.  The statement of the lemma is in
\isa{sat\_hyp\_fuel} and \isa{sat\_fuel}, and the bridges of \S\ref{sec:satat}
convert at the boundary.  \isa{sat\_hyp\_atD} runs downward at the two points
where the induction hands control to \isa{base} and \isa{step}, both of which are
stated in the unindexed predicates.  \isa{sat\_hyp\_at\_exI} runs upward twice,
once in the tail case where the induction hypothesis needs a budget for the
hypotheses of another judgment, and once at the end where the lemma's own
assumption has to produce one.

Instantiating the principle and discharging the step case with
\isa{valid\_step\_sound\_fuel} gives soundness for every judgment occurring in a
checked list.

\thy{bga-check-list-sound}

\noindent
Finally \isa{is\_valid\_proof} is unfolded.  Its two guards are decided with
\isa{cases\_bool}, the head of the proof is shown to be a member of it by
\isa{cons\_reconstr} and \isa{mem\_cons\_head}, and the previous lemma applies.

\thy{bga-bridge}

\subsection{Consistency}
\label{sec:consistency}

The two formula constructors are defined, and their satisfaction is unpacked.

\thy{bga-mk-eq}
\thy{bga-sat-eqE}
\thy{bga-sat-neqE}

\noindent
Both unpacking lemmas go through \isa{tag\_pack\_F} and \isa{load\_pack\_F} to
recover the two sides from the code, and then through the definition of
\isa{sat\_fuel}.  Groundedness of proof checking, the last assumption of the
skeleton, is \isa{proof\_is\_bool}, obtained from \isa{check\_list\_bool} by
list induction and from a chain of \isa{condTB} steps over the branches of
\isa{valid\_step}.

With these in place, \isa{bga\_fuller}, the union of the step-indexed semantics
with the proof checker, is declared a sublocale of \isa{consistent}.

\thy{bga-sublocale}

\noindent
Each obligation is discharged by one of the lemmas above:
\isa{mk\_eq\_N'} and \isa{mk\_neq\_N'} for habeas quid,
\isa{proof\_is\_bool} for groundedness of the checker,
\isa{sat\_hyp\_fuel\_nil} for the empty hypothesis list,
\isa{evals\_functional} for determinism,
\isa{sat\_fuel\_mk\_eqE} and \isa{sat\_fuel\_mk\_neqE} for the two unpacking
assumptions, and \isa{soundness\_bridge\_fuel} for soundness.  The theorem
\isa{syntactically\_consistent} of \S\ref{sec:skeleton} becomes available
inside \isa{bga\_fuller}, and it remains to exhibit an encoding that satisfies
the interface.

\section{The concrete encoding}
\label{sec:encode}

\texttt{encode.thy} instantiates the interface of \S\ref{sec:encoder-interface}
with a G\"odel numbering built from the Cantor pairing of \S\ref{sec:cpair},
supplies concrete recursive definitions for every function the checker is built
from, discharges all the assumptions, and instantiates the consistency theorem.

\subsection{The encoders}

A code is a pair whose first component is the constructor tag and whose second
is the payload.  Formulas use the pairing directly, and terms compose it with
an involution that exchanges $0$ and $1$.

\thy{enc-encoders}

\begin{remarkx}[Why \isa{swap01}]
The interface requires \isa{tag\_T 0 = T\_ZERO}, so the code $0$ must denote the
term $0$.  With \isa{T\_ZERO = 1} and \isa{cpx 0 = 0}, taking
\isa{tag\_T = cpx} would make the code $0$ a variable reference, since
\isa{T\_VAR = 0}, and the structural-decrease assumption
\isa{t ≠ 0 ⟹ load\_T t < t = 1} would then have no atomic base to stop at.
Composing with \isa{swap01} on both the packing and the unpacking side relabels
those two tags and leaves everything else alone, and
\isa{swap01 (swap01 t) = t} makes the composition cancel.  Formulas need no
adjustment, since \isa{F\_EQ = 0} already.
\end{remarkx}

\thy{enc-swap01}

\noindent
\isa{swap01\_N} and \isa{swap01\_inv} are proved by \isa{cases bool:} on the
guards.  The habeas quid premise \isa{t N} is what makes \isa{t = 0} grounded
and the case split available.

\subsection{Concrete definitions}

Every function the locales left abstract is introduced in a single
\isa{axiomatization} block whose axioms are their recursive \isa{:=} equations.
This is the definitional mechanism of \S\ref{sec:defmech} applied at scale, and
it is why a self-referential evaluator can be written down at all: no
termination obligation is incurred when the definition is made.

\thy{enc-signature}

\noindent
Thirty-three constants: the two evaluators, the three substitutions, the
freshness predicates, the assignment update, and the checker hierarchy down to
\isa{is\_valid\_proof}.  Their equations repeat the locale assumptions
verbatim, for instance the direct evaluator

\thy{enc-eval}

\noindent
and the assignment update

\thy{enc-asn-put}

\noindent
The background definition list stays opaque.  Only its termination is assumed:

\begin{center}
\isa{axiomatization dfns :: "dfn" where dfns\_N: "dfns N"}
\end{center}

\noindent
so the result holds for every terminating definition list, including lists
whose definitions diverge on all arguments.

\subsection{Discharging the assumptions}

What remains is the arithmetic that the locale structure confined to this one
file.

\paragraph{Habeas quid.}
Each of the six encoder operations preserves termination, by the corresponding
pairing lemma.

\thy{enc-hq}

\paragraph{Round-trips.}
Injectivity is the projection equations of \S\ref{sec:cpair} composed with
\isa{swap01\_inv}.  The retraction property is \isa{cpair\_reconstr}, that a
number is the pair of its projections.

\thy{enc-roundtrip}

\paragraph{Structural decrease.}
Both payload extractors are \isa{cpy}, so the assumption reduces to strict
decrease of the second projection on nonzero arguments, which is already
available in \texttt{GD.thy}.  The base case is
\isa{swap01 (cpx 0) = swap01 0 = 1}.

\thy{enc-decrease}

\paragraph{Monotonicity.}
This is the assumption that needs work.  The chain begins with
\isa{pair\_step}, that a single successor step in the second component does not
decrease the pair,
\begin{center}
\isa{pair\_step: ⟦a N; k N⟧ ⟹ ⟨a, k⟩ ≤ ⟨a, S k⟩ = 1},
\end{center}
which follows from \isa{cpair\_suc} and repeated use of
\isa{leq\_monotone\_add\_r}.  \isa{pair\_mono\_2} closes it under induction on
the second component, and the two interface obligations follow.

\thy{enc-mono}

\noindent
The induction in \isa{pair\_mono\_2} is a good miniature of the style of the
whole development.  The statement being inducted on is the object-level
implication \isa{(x ≤ y = 1) ⟶ (⟨a, x⟩ ≤ ⟨a, y⟩ = 1)}, every \isa{implI} step
first discharges a groundedness obligation through \isa{eqBool}, and the case
analysis in the step case goes through \isa{cases\_bool}.

\subsection{Instantiation and the theorem}

With the assumptions available the locale is interpreted.  The interpretation
supplies, in order, the six encoder operations, the three freshness predicates,
the definition list and its membership test, the step-indexed evaluator, the
three substitutions and the assignment update, and the checker hierarchy.

\thy{enc-interpretation}

\noindent
Every goal generated by \isa{unfold\_locales} is closed by \isa{fact}: the
structural obligations by the lemmas above, and the definitional obligations by
the \isa{:=} axioms of the \isa{axiomatization} block, which are literally the
equations the locale assumes.

\thy{enc-theorem}

\noindent
Read out, the theorem says that for all terminating $a$, $b$, $p_1$ and $p_2$,
it is not the case that $p_1$ codes a \BGA{} proof of $\enc{a = b}$ from no
hypotheses while $p_2$ codes a \BGA{} proof of $\enc{a \neq b}$ from no
hypotheses.  Every symbol in that statement, the quantifiers, the negation, the
conjunction and the proof checker, belongs to \QGA.

\section{Reduction of \QGA{} to \BGA}
\label{sec:reduction}

\S\ref{sec:proof} proves \BGA{} consistent inside \QGA.  To turn that into a
self-consistency result, \QGA{} itself has to be reduced to \BGA, so that the
system doing the proving is no stronger than the system being proved
consistent.  This section is to be written.


\section{Status}

The three theories are complete in the sense that \texttt{encode.thy}
discharges every assumption of \isa{bga\_fuller} by \isa{unfold\_locales} and
derives \isa{BGA\_syntactically\_consistent} from
\isa{conc2.syntactically\_consistent}.  What remains for this document is the
detail of the soundness proofs summarized in \S\ref{sec:proof}, and the content
of \S\ref{sec:reduction}.

Three points about the development itself are worth recording.

\begin{enumerate}
\item The concrete functions of \texttt{encode.thy}, the evaluator, the
      substitutions and the checkers, are introduced by an
      \isa{axiomatization} block whose axioms are their recursive \isa{:=}
      equations.  This is the native definitional mechanism of \GA{}
      (\S\ref{sec:defmech}), which is sound because \GA{} places no termination
      requirement on a recursive definition.  It does mean the functions are
      postulated to satisfy their equations rather than built by a conservative
      definitional package.
\item The exact evaluator \isa{eval} of \isa{bga\_semantics} is retained in the
      sources but is not on the critical path.  The consistency argument runs
      through the step-indexed \isa{eval\_fuel} and the relation \isa{evals}
      derived from it (\S\ref{sec:fuel}), for which the determinism lemma is
      available.
\item An earlier version of the development, the locale \isa{bga\_full} and its
      interpretation, is retained inside a comment block in
      \texttt{BGA\_on\_GA.thy} and is superseded by \isa{bga\_fuller}.  Nothing
      described here depends on it.
\end{enumerate}

\end{document}
